\documentclass[10pt,border=1pt]{standalone}
\usepackage{amsmath,amssymb,booktabs,array}
\usepackage{atlas-figs}
\begin{document}%
\begin{tabular}{@{}c@{}}\begin{tikzpicture}[bf,x=1cm,y=1cm]
\draw[black!25,line width=.3pt] (0.000,.4) -- (0.000,-10.61);\node[font=\footnotesize,anchor=south] at (0.000,.4) {$8$};
\draw[black!25,line width=.3pt] (0.957,.4) -- (0.957,-10.61);\node[font=\footnotesize,anchor=south] at (0.957,.4) {$12$};
\draw[black!25,line width=.3pt] (1.780,.4) -- (1.780,-10.61);\node[font=\footnotesize,anchor=south] at (1.780,.4) {$17$};
\draw[black!25,line width=.3pt] (2.594,.4) -- (2.594,-10.61);\node[font=\footnotesize,anchor=south] at (2.594,.4) {$24$};
\draw[black!25,line width=.3pt] (3.346,.4) -- (3.346,-10.61);\node[font=\footnotesize,anchor=south] at (3.346,.4) {$33$};
\draw[black!25,line width=.3pt] (4.327,.4) -- (4.327,-10.61);\node[font=\footnotesize,anchor=south] at (4.327,.4) {$50$};
\draw[black!25,line width=.3pt] (5.284,.4) -- (5.284,-10.61);\node[font=\footnotesize,anchor=south] at (5.284,.4) {$75$};
\draw[black!25,line width=.3pt] (5.963,.4) -- (5.963,-10.61);\node[font=\footnotesize,anchor=south] at (5.963,.4) {$100$};
\draw[black!25,line width=.3pt] (6.921,.4) -- (6.921,-10.61);\node[font=\footnotesize,anchor=south] at (6.921,.4) {$150$};
\draw[black!25,line width=.3pt] (7.600,.4) -- (7.600,-10.61);\node[font=\footnotesize,anchor=south] at (7.600,.4) {$200$};
\node[font=\footnotesize,anchor=south west,inner sep=1pt] at (0,.95) {number of vertices of $T$};
\node[font=\footnotesize\bfseries,anchor=east,inner sep=1pt] at (-.2,.2) {open, and tested up to};
\draw[line width=1.5pt,dash pattern=on 3pt off 2pt] (2.388,-0.45999999999999996) -- (5.963,-0.45999999999999996);\node[font=\footnotesize,fill=white,inner sep=1pt,anchor=west] at (6.013,-0.45999999999999996) {$100$};
\draw[line width=.7pt,fill=white] (0,-0.57) rectangle (2.388,-0.35);
\fill[black!80] (0,-0.57) rectangle (2.388,-0.35);
\node[font=\footnotesize,fill=white,inner sep=1pt,anchor=south west] at (2.428,-0.37999999999999995) {$22$};
\node[font=\footnotesize,anchor=east,align=right,inner sep=1pt] at (-.2,-0.45999999999999996) {at least $2^m$ good states, $m$ the fewest\\cycles of a cover: every root};
\draw[line width=.7pt,fill=white] (0,-1.49) rectangle (2.594,-1.2699999999999998);
\fill[black!80] (0,-1.49) rectangle (1.637,-1.2699999999999998);
\node[font=\footnotesize,fill=white,inner sep=1pt,anchor=south west] at (2.634,-1.2999999999999998) {$16, 24$};
\node[font=\footnotesize,anchor=east,align=right,inner sep=1pt] at (-.2,-1.38) {the same bound, on triangulations\\without separating triangles};
\draw[line width=1.5pt,dash pattern=on 3pt off 2pt] (2.493,-2.3) -- (7.600,-2.3);\node[font=\footnotesize,fill=white,inner sep=1pt,anchor=west] at (7.650,-2.3) {$200$};
\draw[line width=.7pt,fill=white] (0,-2.4099999999999997) rectangle (2.493,-2.19);
\fill[black!80] (0,-2.4099999999999997) rectangle (2.163,-2.19);
\node[font=\footnotesize,fill=white,inner sep=1pt,anchor=south west] at (2.533,-2.2199999999999998) {$20, 23$};
\node[font=\footnotesize,anchor=east,align=right,inner sep=1pt] at (-.2,-2.3) {pair descent that never adds a cycle\\reaches a good pair};
\draw[line width=1.5pt,dash pattern=on 3pt off 2pt] (1.780,-3.2199999999999998) -- (6.394,-3.2199999999999998);\node[font=\footnotesize,fill=white,inner sep=1pt,anchor=west] at (6.444,-3.2199999999999998) {$120$};
\draw[line width=.7pt,fill=white] (0,-3.3299999999999996) rectangle (1.780,-3.11);
\fill[black!80] (0,-3.3299999999999996) rectangle (0.957,-3.11);
\node[font=\footnotesize,fill=white,inner sep=1pt,anchor=south west] at (1.820,-3.1399999999999997) {$12, 17$};
\node[font=\footnotesize,anchor=east,align=right,inner sep=1pt] at (-.2,-3.2199999999999998) {exchange descent on rainbow trees\\reaches a good state};
\draw[line width=.7pt,fill=white] (0,-4.25) rectangle (2.042,-4.029999999999999);
\fill[black!80] (0,-4.25) rectangle (2.042,-4.029999999999999);
\node[font=\footnotesize,fill=white,inner sep=1pt,anchor=south west] at (2.082,-4.06) {$19$};
\node[font=\footnotesize,anchor=east,align=right,inner sep=1pt] at (-.2,-4.14) {without separating triangles,\\some contraction keeps a pair};
\draw[line width=.7pt,fill=white] (0,-5.17) rectangle (1.915,-4.949999999999999);
\fill[black!80] (0,-5.17) rectangle (1.915,-4.949999999999999);
\node[font=\footnotesize,fill=white,inner sep=1pt,anchor=south west] at (1.955,-4.9799999999999995) {$18$};
\node[font=\footnotesize,anchor=east,align=right,inner sep=1pt] at (-.2,-5.06) {a contraction never has more\\good states than the original};
\node[font=\footnotesize\bfseries,anchor=east,inner sep=1pt] at (-.2,-6.02) {refuted, counterexample known at};
\draw[line width=.4pt,black!55] (0,-6.489999999999999) -- (0.957,-6.489999999999999);\draw[line width=1.6pt] (0.857,-6.589999999999999) -- (1.057,-6.39) (0.857,-6.39) -- (1.057,-6.589999999999999);
\node[font=\footnotesize,inner sep=1pt,anchor=west] at (1.207,-6.489999999999999) {$12$};
\node[font=\footnotesize,anchor=east,align=right,inner sep=1pt] at (-.2,-6.489999999999999) {the highest state with a partner is good};
\draw[line width=.4pt,black!55] (0,-6.989999999999999) -- (1.915,-6.989999999999999);\draw[line width=1.6pt] (1.815,-7.089999999999999) -- (2.015,-6.89) (1.815,-6.89) -- (2.015,-7.089999999999999);
\draw[line width=.8pt] (1.915,-6.989999999999999) circle[radius=.19];
\node[font=\footnotesize,inner sep=1pt,anchor=west] at (2.165,-6.989999999999999) {$18$};
\node[font=\footnotesize,anchor=east,align=right,inner sep=1pt] at (-.2,-6.989999999999999) {an extreme cover has a good orientation};
\draw[line width=.4pt,black!55] (0,-7.659999999999999) -- (1.915,-7.659999999999999);\draw[line width=1.6pt] (1.815,-7.759999999999999) -- (2.015,-7.56) (1.815,-7.56) -- (2.015,-7.759999999999999);
\draw[line width=.8pt] (1.915,-7.659999999999999) circle[radius=.19];
\node[font=\footnotesize,inner sep=1pt,anchor=west] at (2.165,-7.659999999999999) {$18$};
\node[font=\footnotesize,anchor=east,align=right,inner sep=1pt] at (-.2,-7.659999999999999) {some good state of a contraction\\lifts directly};
\draw[line width=.4pt,black!55] (0,-8.33) -- (2.594,-8.33);\draw[line width=1.6pt] (2.494,-8.43) -- (2.694,-8.23) (2.494,-8.23) -- (2.694,-8.43);
\node[font=\footnotesize,inner sep=1pt,anchor=west] at (2.844,-8.33) {$24$};
\node[font=\footnotesize,anchor=east,align=right,inner sep=1pt] at (-.2,-8.33) {a cover one step from an extreme has one};
\draw[line width=.4pt,black!55] (0,-8.83) -- (3.346,-8.83);\draw[line width=1.6pt] (3.246,-8.93) -- (3.446,-8.73) (3.246,-8.73) -- (3.446,-8.93);
\node[font=\footnotesize,inner sep=1pt,anchor=west] at (3.596,-8.83) {$33$};
\node[font=\footnotesize,anchor=east,align=right,inner sep=1pt] at (-.2,-8.83) {an extreme or irreducible cover has one};
\draw[line width=.4pt,black!55] (0,-9.33) -- (3.346,-9.33);\draw[line width=1.6pt] (3.246,-9.43) -- (3.446,-9.23) (3.246,-9.23) -- (3.446,-9.43);
\node[font=\footnotesize,inner sep=1pt,anchor=west] at (3.596,-9.33) {$33$};
\node[font=\footnotesize,anchor=east,align=right,inner sep=1pt] at (-.2,-9.33) {every maximal chain of covers has one};
\draw[line width=.4pt,black!55] (0,-9.83) -- (3.346,-9.83);\draw[line width=1.6pt] (3.246,-9.93) -- (3.446,-9.73) (3.246,-9.73) -- (3.446,-9.93);
\node[font=\footnotesize,inner sep=1pt,anchor=west] at (3.596,-9.83) {$33$};
\node[font=\footnotesize,anchor=east,align=right,inner sep=1pt] at (-.2,-9.83) {at least half of all covers have one};
\draw[line width=.4pt,black!55] (0,-10.33) -- (3.679,-10.33);\draw[line width=1.6pt] (3.579,-10.43) -- (3.779,-10.23) (3.579,-10.23) -- (3.779,-10.43);
\node[font=\footnotesize,inner sep=1pt,anchor=west] at (3.929,-10.33) {$38$};
\node[font=\footnotesize,anchor=east,align=right,inner sep=1pt] at (-.2,-10.33) {the evict walk never repeats};
\fill[black!80] (-5.8,-11.19) rectangle (-5.0,-10.97);\node[right,font=\footnotesize] at (-5.0,-11.08) {every case, two programs};
\draw[line width=.7pt,fill=white] (-1.2999999999999998,-11.19) rectangle (-0.5,-10.97);\node[right,font=\footnotesize] at (-0.5,-11.08) {every case, one program};
\draw[line width=1.5pt,dash pattern=on 3pt off 2pt] (3.2,-11.08) -- (4.000000000000001,-11.08) node[right,font=\footnotesize] {random samples};
\draw[line width=1.6pt] (6.6000000000000005,-11.18) -- (6.8,-10.98) (6.6000000000000005,-10.98) -- (6.8,-11.18);\draw[line width=.8pt] (6.7,-11.08) circle[radius=.19];\node[right,font=\footnotesize] at (6.8999999999999995,-11.08) {smallest};
\end{tikzpicture}\\[1mm]\begin{minipage}{14cm}\centering\footnotesize (a) How far each open statement was tested, and where a counterexample to each refuted one is known. Contractions are at nonroot degree-four vertices. Logarithmic scale.\end{minipage}\\\noalign{\vskip 5mm}\begin{tabular}{@{}c@{\hspace{5mm}}c@{}}\begin{tikzpicture}[bf,bf small,x=0.47cm,y=0.47cm]
\path[use as bounding box] (-5.80,-5.80) rectangle (5.80,5.80);
\fill[bfblack,even odd rule] (0,0) circle[radius=5.55] (0,0) circle[radius=5.0];
\fill[bfblack] (-4.330,-2.500) -- (-2.133,-0.405) -- (-2.200,3.528) -- (-4.330,-2.500) -- cycle;
\fill[bfblack] (-4.330,-2.500) -- (-0.344,-4.394) -- (-1.614,-2.366) -- (-4.330,-2.500) -- cycle;
\fill[bfblack] (-2.133,-0.405) -- (0.764,-0.167) -- (-0.963,1.243) -- (-2.133,-0.405) -- cycle;
\fill[bfblack] (-2.133,-0.405) -- (-1.614,-2.366) -- (0.900,-2.302) -- (-2.133,-0.405) -- cycle;
\fill[bfblack] (-0.344,-4.394) -- (4.330,-2.500) -- (0.900,-2.302) -- (-0.344,-4.394) -- cycle;
\fill[bfblack] (4.330,-2.500) -- (2.799,2.992) -- (2.506,-0.651) -- (4.330,-2.500) -- cycle;
\fill[bfblack] (0.000,5.000) -- (-2.200,3.528) -- (2.799,2.992) -- (0.000,5.000) -- cycle;
\fill[bfblack] (-2.200,3.528) -- (-0.963,1.243) -- (0.610,2.065) -- (-2.200,3.528) -- cycle;
\fill[bfblack] (0.764,-0.167) -- (2.799,2.992) -- (0.610,2.065) -- (0.764,-0.167) -- cycle;
\fill[bfblack] (0.764,-0.167) -- (0.900,-2.302) -- (2.506,-0.651) -- (0.764,-0.167) -- cycle;
\draw[bf edge] (0,0) circle[radius=5.0];
\draw[bf edge] (-4.330,-2.500) -- (-2.133,-0.405);
\draw[bf edge] (-4.330,-2.500) -- (-1.614,-2.366);
\draw[bf edge] (-4.330,-2.500) -- (-0.344,-4.394);
\draw[bf edge] (-4.330,-2.500) -- (-2.200,3.528);
\draw[bf edge] (-2.133,-0.405) -- (-1.614,-2.366);
\draw[bf edge] (-2.133,-0.405) -- (-2.200,3.528);
\draw[bf edge] (-2.133,-0.405) -- (-0.963,1.243);
\draw[bf edge] (-2.133,-0.405) -- (0.764,-0.167);
\draw[bf edge] (-2.133,-0.405) -- (0.900,-2.302);
\draw[bf edge] (-1.614,-2.366) -- (-0.344,-4.394);
\draw[bf edge] (-1.614,-2.366) -- (0.900,-2.302);
\draw[bf edge] (-0.344,-4.394) -- (4.330,-2.500);
\draw[bf edge] (-0.344,-4.394) -- (0.900,-2.302);
\draw[bf edge] (4.330,-2.500) -- (0.900,-2.302);
\draw[bf edge] (4.330,-2.500) -- (2.506,-0.651);
\draw[bf edge] (4.330,-2.500) -- (2.799,2.992);
\draw[bf edge] (0.000,5.000) -- (-2.200,3.528);
\draw[bf edge] (0.000,5.000) -- (2.799,2.992);
\draw[bf edge] (-2.200,3.528) -- (-0.963,1.243);
\draw[bf edge] (-2.200,3.528) -- (2.799,2.992);
\draw[bf edge] (-2.200,3.528) -- (0.610,2.065);
\draw[bf edge] (-0.963,1.243) -- (0.764,-0.167);
\draw[bf edge] (-0.963,1.243) -- (0.610,2.065);
\draw[bf edge] (0.764,-0.167) -- (0.900,-2.302);
\draw[bf edge] (0.764,-0.167) -- (2.506,-0.651);
\draw[bf edge] (0.764,-0.167) -- (2.799,2.992);
\draw[bf edge] (0.764,-0.167) -- (0.610,2.065);
\draw[bf edge] (0.900,-2.302) -- (2.506,-0.651);
\draw[bf edge] (2.506,-0.651) -- (2.799,2.992);
\draw[bf edge] (2.799,2.992) -- (0.610,2.065);
\fill[af tint] (-2.133,-0.405) -- (0.900,-2.302) -- (0.764,-0.167) -- (-2.133,-0.405) -- cycle;
\begin{scope}\clip (-2.133,-0.405) -- (0.764,-0.167) -- (-0.963,1.243) -- (-2.133,-0.405) -- cycle;\path[bf s] (-0.963,1.243) -- ++(-125.37:0.709) arc[start angle=-125.37,end angle=-39.23,radius=0.709] -- cycle;\end{scope}
\begin{scope}\clip (-2.133,-0.405) -- (-1.614,-2.366) -- (0.900,-2.302) -- (-2.133,-0.405) -- cycle;\path[bf s] (-1.614,-2.366) -- ++(1.46:0.648) arc[start angle=1.46,end angle=104.82,radius=0.648] -- cycle;\end{scope}
\begin{scope}\clip (0.764,-0.167) -- (0.900,-2.302) -- (2.506,-0.651) -- (0.764,-0.167) -- cycle;\path[bf s] (2.506,-0.651) -- ++(164.47:0.723) arc[start angle=164.47,end angle=225.79,radius=0.723] -- cycle;\end{scope}
\begin{scope}\clip (-4.330,-2.500) -- (-2.133,-0.405) -- (-2.200,3.528) -- (-4.330,-2.500) -- cycle;\path[fill=black!45,draw=black!45] (-2.133,-0.405) -- ++(90.98:0.572) arc[start angle=90.98,end angle=223.64,radius=0.572] -- cycle;\end{scope}
\begin{scope}\clip (-4.330,-2.500) -- (-0.344,-4.394) -- (-1.614,-2.366) -- (-4.330,-2.500) -- cycle;\path[fill=black!45,draw=black!45] (-0.344,-4.394) -- ++(122.06:0.957) arc[start angle=122.06,end angle=154.58,radius=0.957] -- cycle;\end{scope}
\begin{scope}\clip (-0.344,-4.394) -- (4.330,-2.500) -- (0.900,-2.302) -- (-0.344,-4.394) -- cycle;\path[fill=black!45,draw=black!45] (0.900,-2.302) -- ++(-120.74:0.608) arc[start angle=-120.74,end angle=-3.30,radius=0.608] -- cycle;\end{scope}
\begin{scope}\clip (4.330,-2.500) -- (2.799,2.992) -- (2.506,-0.651) -- (4.330,-2.500) -- cycle;\path[fill=black!45,draw=black!45] (2.799,2.992) -- ++(-94.60:1.462) arc[start angle=-94.60,end angle=-74.42,radius=1.462] -- cycle;\end{scope}
\begin{scope}\clip (0.000,5.000) -- (-2.200,3.528) -- (2.799,2.992) -- (0.000,5.000) -- cycle;\path[fill=black!45,draw=black!45] (-2.200,3.528) -- ++(-6.12:1.042) arc[start angle=-6.12,end angle=33.79,radius=1.042] -- cycle;\end{scope}
\begin{scope}\clip (-2.200,3.528) -- (-0.963,1.243) -- (0.610,2.065) -- (-2.200,3.528) -- cycle;\path[fill=black!45,draw=black!45] (0.610,2.065) -- ++(152.50:0.710) arc[start angle=152.50,end angle=207.59,radius=0.710] -- cycle;\end{scope}
\begin{scope}\clip (0.764,-0.167) -- (2.799,2.992) -- (0.610,2.065) -- (0.764,-0.167) -- cycle;\path[fill=black!45,draw=black!45] (0.764,-0.167) -- ++(57.21:0.895) arc[start angle=57.21,end angle=93.95,radius=0.895] -- cycle;\end{scope}
\draw[bf Q,line width=2pt] (-2.133,-0.405) -- (0.764,-0.167);
\draw[bf Q,line width=2pt] (-2.133,-0.405) -- (0.900,-2.302);
\draw[bf Q,line width=2pt] (0.764,-0.167) -- (0.900,-2.302);
\draw[bf root] (-4.330,-2.500) circle[radius=1.4pt];
\fill[bf vertex] (-2.133,-0.405) circle[radius=.8pt];
\fill[bf vertex] (-1.614,-2.366) circle[radius=.8pt];
\fill[bf vertex] (-0.344,-4.394) circle[radius=.8pt];
\draw[bf root] (4.330,-2.500) circle[radius=1.4pt];
\draw[bf root] (0.000,5.000) circle[radius=1.4pt];
\fill[bf vertex] (-2.200,3.528) circle[radius=.8pt];
\fill[bf vertex] (-0.963,1.243) circle[radius=.8pt];
\fill[bf vertex] (0.764,-0.167) circle[radius=.8pt];
\fill[bf vertex] (0.900,-2.302) circle[radius=.8pt];
\fill[bf vertex] (2.506,-0.651) circle[radius=.8pt];
\fill[bf vertex] (2.799,2.992) circle[radius=.8pt];
\fill[bf vertex] (0.610,2.065) circle[radius=.8pt];
\end{tikzpicture} & \begin{minipage}[b][5.4cm][c]{7.6cm}\footnotesize (b) The mechanism behind the failures at the extremes, on the state $s_{9}$ of the running example. The three black faces round the hatched white face all put their wedge (blue) at the corner away from it, so their opposite edges close a triangle (heavy dashes). The 12-vertex counterexample on page 16 fails through a face of this kind.\end{minipage}\end{tabular}\end{tabular}%
\end{document}
